\documentclass{article}
\usepackage[final,nonatbib]{neurips_2019}
\usepackage[utf8]{inputenc}
\usepackage[T1]{fontenc}
\usepackage{amsmath,booktabs,microtype}
\usepackage[hidelinks]{hyperref}
\newcommand{\ZoneTest}{22.11}
\newcommand{\PooledTest}{21.06}
\newcommand{\MlpTest}{21.19}
\newcommand{\GatTest}{20.97}
\newcommand{\GruTest}{20.98}
\newcommand{\ValPaired}{428,213}
\newcommand{\TestPaired}{633,885}
\newcommand{\ValAll}{428,427}
\newcommand{\TestAll}{634,578}
\newcommand{\MlpGatDifference}{0.21}
\newcommand{\SelfGatDifference}{0.09}
\newcommand{\GatRmse}{34.92}
\newcommand{\GruRmse}{47.02}
\newcommand{\GruBestEpoch}{29}
\newcommand{\GruEpochs}{37}
\newcommand{\MlpCiLower}{0.20}
\newcommand{\MlpCiUpper}{0.86}
\newcommand{\SelfCiLower}{0.13}
\newcommand{\SelfCiUpper}{0.46}

\title{Weather, Interconnections, and Energy Prices\\
  \small\normalfont Stanford CS224W Project Milestone v3 --- Draft, November 9, 2026}
\author{Project team\\\small Author details to be completed}
\begin{document}
\maketitle
\begin{abstract}
Across 49 European market identities, we compare persistence, independent/pooled
XGBoost, static/GRU MLP and self-loop/real GAT. Pooled XGBoost, static MLP and real
GAT achieve \PooledTest, \MlpTest\ and \GatTest\ EUR/MWh matched test MAE;
the verified residual GRU achieves \GruTest\ in one seed. Graph gains are modest
after controlling for pooled learning. A numerical audit invalidates the earlier
negative GRU interpretation. These are exploratory retrospective benchmarks with
unresolved source vintages.
\end{abstract}

\section{Approach}
At 12:00 Europe/Brussels on calendar day $D-1$, predict every hourly day-ahead
price of delivery day $D$, including 23/25-hour daylight-saving days. Weekly
persistence copies the same market's price 168 elapsed hours earlier.
Independent and pooled XGBoost \cite{chen} distinguish shared learning from
cross-market graph messages. A pooled MLP provides a classical neural reference.
Self-loop GATv2 \cite{gat} retains the architecture without cross-market messages;
real GAT adds physical pairs observed during training. These GATs share features
and training settings. Node-lifetime masks preserve historical geography.
JAO constraints enter as flat summaries, not extra graph nodes.

\section{Experiments}
\paragraph{Data and partitions.}
The canonical master contains 5,887,954 zone-hours over 62 historical price-area
identities, January 2016--October 2026. ENTSO-E prices/load/flows use pinned v1.2
releases; other sources are an ERA5-lineage weather export and original CWE/CORE
JAO observations. Weather's requested v1.1 release identity is unverified.
Missing values, negative prices and separate DE\_AT\_LU/DE\_LU identities are
preserved; 49 identities have labels in all required partitions. Targets are
never filled or filtered for feature completeness. Fit uses January 2016--March
2024, validation April 2024--March 2025, and test April 2025--September 2026 by
Brussels delivery date. October stays unassigned.

\paragraph{Static controls.}
All five static learned controls use identical richer inputs: local
hour/weekday/month, guarded 24/168-hour prices, load forecast, four lag-48 weather
variables and three lag-48 JAO summaries (constraint count, minimum RAM ratio,
domain presence). Trees retain NaNs; neural imputation/scaling fits training
features only and includes missingness indicators. Models fit
$z=\operatorname{asinh}(p/50)$ with L1 loss and invert with $50\sinh(\hat z)$.
Trees use depth six, at most 600 rounds and 50-round validation patience.
MLP has two ELU hidden layers matched by parameter count to GAT; GAT has two
layers, 64 channels and four heads. Both use Adam 0.001, dropout 0.1, at most
30 epochs and patience five. All five controls use seeds 224--228.

\paragraph{Partner-inspired temporal model.}
The GRU--MLP adapts Sasha's recipe on our canonical targets, rather than importing
his separate 43-zone score. A 64-unit GRU encodes 168 completed hours of price
lags, load, derived net inflow and four weather channels. Two 128-unit ReLU MLP
layers combine its state with 11 day inputs. The target is
$r=z-\operatorname{asinh}(b/50)$, with weekly-price base $b$, else guarded daily
price, else zero; this fallback reconstructs predictions only. Residual L1 is
augmented by 0.5 times connected-zone residual-contrast L1. Day weather uses
lag 48. Seed 0, 16-day batches, Adam 0.001 and patience eight select epoch
\GruBestEpoch\ of \GruEpochs. This is a different recipe and one seed,
not a controlled GRU or weather ablation.

\paragraph{Metrics.}
MAE/RMSE are EUR/MWh after inversion. Relative MAE divides model MAE by
persistence MAE on identical hours: 0.72 means 28\% less error, not 28\% of the
price. Table~\ref{tab:main} uses \ValPaired\ validation and \TestPaired\
test labels with weekly history. Full scoring retains all \ValAll\ validation
and \TestAll\ test labels, including unavailable histories. Validation selects
checkpoints; no best test seed is chosen.

\begin{table}[!ht]
\centering\small
\begin{tabular}{lrrrr}
\toprule
Forecasting method & Seeds & Val. MAE & Test MAE & Test rMAE \\
\midrule
Weekly persistence & -- & 30.18 & 29.13 & 1.0000 \\
\midrule
Per-zone XGBoost & 5 & 21.23 $\pm$ 0.05 & 22.11 $\pm$ 0.03 & 0.7589 \\
Pooled XGBoost & 5 & 20.74 $\pm$ 0.01 & 21.06 $\pm$ 0.01 & 0.7228 \\
\midrule
Static pooled MLP & 5 & 21.44 $\pm$ 0.23 & 21.19 $\pm$ 0.21 & 0.7272 \\
GRU + MLP (partner)$^\dagger$ & 1 & \textbf{20.32} & 20.98 & 0.7202 \\
\midrule
Self-loop GAT & 5 & 21.27 $\pm$ 0.09 & 21.06 $\pm$ 0.08 & 0.7230 \\
Real-graph GAT & 5 & 20.98 $\pm$ 0.11 & \textbf{20.97 $\pm$ 0.13} & \textbf{0.7199} \\
\bottomrule
\end{tabular}

\caption{Same 49-market cohort and matched hours. Static learned rows: five-seed
mean $\pm$ sample SD, not ensembles. $\dagger$: one verified seed, different
inputs/target/loss. Bold denotes the smallest reported mean, not significance
or a controlled GRU advantage. Lower is better.}
\label{tab:main}
\end{table}

\paragraph{Results.}
Pooling reduces XGBoost test MAE from \ZoneTest\ to \PooledTest.
Real GAT improves test MAE by \MlpGatDifference\ versus static MLP and
\SelfGatDifference\ versus self-loop GAT. Corresponding validation-week
improvement intervals are [\MlpCiLower, \MlpCiUpper] and
[\SelfCiLower, \SelfCiUpper] EUR/MWh (2,000 shared-week resamples, conditional
on fitted models). Pooled XGBoost has better validation MAE than real GAT;
their small test gap does not justify retrospective model choice.
The GRU has similar matched MAE but all-target test RMSE \GruRmse{} versus
\GatRmse\ for real GAT. Gains vary by market; graph superiority is not universal.

\paragraph{Numerical audit.}
The original controlled GRU--MLP/GAT runs are omitted: all ten checkpoints
disagree with CPU on sampled shuffled training-shaped inputs under fused ROCm,
despite replaying saved consecutive inference batches. Native GPU agrees with
CPU; the negative temporal conclusion is withdrawn. The accepted partner-inspired
run uses native GPU operations and passes sampled CPU float64, checkpoint-replay,
target-signature and saved-score checks. This is not an exhaustive backend proof.

\paragraph{Availability and limitations.}
ERA5 is retrospective reanalysis with roughly five-day release latency
\cite{era5}, incompatible with recent lag-48 inputs. Load/JAO/price revision
vintages remain unresolved. Exact timestamp joins and autumn-DST guards prevent
lag24 from reaching the target delivery day, but cannot reconstruct publication
records. Scores are conditional, not certified operational forecasts. Test data
were previously explored; intervals are not confirmatory evidence.

\section{Future work}
Retrain our original temporal controls with native GPU operations; repeat the
residual GRU across seeds and ablate its residual target and contrast loss.
Build a daily service that records inputs, versions and predictions before
Brussels noon, then scores auction outcomes. Start with persistence and
price-only pooled XGBoost; monitor rolling rMAE, coverage, deadline failures and
fallbacks. Add timely weather observations/archived forecasts and graph challengers
after establishing this prospective baseline.

\begin{thebibliography}{3}\small
\bibitem{chen} T. Chen and C. Guestrin. XGBoost: A scalable tree boosting system.
\emph{KDD}, 2016. \href{https://arxiv.org/abs/1603.02754}{arXiv:1603.02754}.
\bibitem{gat} S. Brody et al. How attentive are graph attention
networks? \emph{ICLR}, 2022. \href{https://arxiv.org/abs/2105.14491}{arXiv:2105.14491}.
\bibitem{era5} ECMWF.
\href{https://www.ecmwf.int/en/forecasts/datasets/era5-hourly-time-series-data-single-levels-1940-present}{ERA5 hourly time-series product description}.
Accessed November 9, 2026.
\end{thebibliography}
\end{document}
